\documentclass[10pt]{article}
\usepackage[a4paper,margin=1.9cm]{geometry}
\usepackage{amsmath,amssymb}
\usepackage{graphicx}
\graphicspath{{../}{./}}
\usepackage{float}
\usepackage[dvipsnames]{xcolor}
\usepackage{enumitem}
\usepackage{titlesec}
\usepackage[hidelinks]{hyperref}

% --- compact layout ---
\setlist[itemize]{leftmargin=1.35em,itemsep=1pt,topsep=2pt,parsep=0pt}
\setlength{\parindent}{0pt}
\setlength{\parskip}{3pt}
\titlespacing*{\section}{0pt}{6pt}{3pt}
\titlespacing*{\subsection}{0pt}{4pt}{2pt}
\titleformat{\section}{\normalfont\large\bfseries\color{RoyalBlue}}{\thesection.}{0.5em}{}
\titleformat{\subsection}{\normalfont\bfseries}{}{0em}{}

\newcommand{\code}[1]{\texttt{\small #1}}
\newcommand{\file}[1]{\textcolor{RoyalBlue}{\texttt{\small #1}}}

% \docheader{title}{subtitle}
\newcommand{\docheader}[2]{%
  \begin{center}
    {\LARGE\bfseries #1}\\[2pt]
    {\small #2}
  \end{center}
  \vspace{-4pt}\hrule\vspace{6pt}
}

\begin{document}

\docheader{Adora --- GPU Memory: Host-Resident Column Batching}%
{As-built notes \quad$\cdot$\quad Ivan Milic \quad$\cdot$\quad 2026-09-20 \quad$\cdot$\quad session summary}

The batched GPU synth (\code{synth\_cube\_map(..., batch=N)}, see \file{adora\_9\_eos\_integration.pdf})
ran a $768^2$ cube but died on larger ones. Cause: \code{batch} bounded the opacity intermediate but
\emph{not} the input residency. This note fixes that so GPU memory is bounded by \code{batch} regardless
of cube size. Everything runs on the GPU by default (the \code{adora-develop} JAX is the CUDA build;
\code{jax.default\_backend()=="gpu"}).

\subsection*{1. The bug}
The flatten uploaded the \textbf{whole} cube to the GPU before the chunk loop:
{\small\begin{verbatim}
flat = [jnp.asarray(cube[k]).reshape(-1, nz) for k in keys]   # all columns, all keys -> GPU
\end{verbatim}}
\code{batch} then sliced arrays that were \emph{already} on the device, so it capped the per-chunk
opacity ($\sim$\,\code{batch}$\cdot n_\lambda\cdot n_z\cdot4\cdot4\cdot8$ B) but left the full input cube
resident. Rough resident inputs (7 keys, $n_z{=}121$, fp64): $768^2\!\to\!\sim$4\,GB (fits),
$1536^2\!\to\!\sim$16\,GB (does not, alongside compute + output).

\subsection*{2. The fix (\file{synth3d.py}, \code{synth\_cube\_map})}
Keep \code{flat} on the \textbf{host} (numpy); move only the current chunk to the GPU and pull its
intensity straight back:
{\small\begin{verbatim}
shp  = cube["T"].shape                              # read shape without uploading anything
flat = [np.asarray(cube[k]).reshape(-1, nz) for k in keys]      # host (numpy)
...
for s in starts:
    sub = [jnp.asarray(a[s:s+batch]) for a in flat]            # host slice -> GPU (only this chunk)
    chunks.append(np.asarray(fn(adata, waves, dz, *sub)))      # intensity -> host, frees the GPU
I = np.concatenate(chunks, axis=0)
\end{verbatim}}
Peak GPU memory is now \code{batch} columns of inputs $+$ their opacity, \textbf{independent of the total
cube size}. The same warm \code{fn} is reused (one compile; a smaller final chunk recompiles once).
The \code{batch=None} path is unchanged: jit uploads the whole cube and returns an on-device JAX array
(still composes with \code{grad}/\code{jit}) --- host batching is only for large forward runs.

\subsection*{3. Verified}
Batched output is \textbf{bit-identical} to the whole-cube result (\code{max|$\Delta$|=0}) for both the
EOS (\code{pg}) and legacy \code{ne/nH} paths, including an uneven final chunk; \code{batch}$\ge$\code{npix}
falls back to the single call. (FAL-C toy cube, forward-only.)

\subsection*{4. Notes / next}
Host RAM now holds the input cube and the concatenated fp64 output ($\sim$15\,GB at $1536^2$) --- the new
ceiling is host memory, not the GPU. The complementary $\sim2\times$ lever is \textbf{float32}: x64 is
currently hardcoded \code{True} in all modules; running \code{jax\_enable\_x64=False} globally would make
the kernel constants (\code{jnp.eye(4)}, Planck) f32 too, halving memory --- not wired as a switch yet, and
the polarised DELO solve / optical-depth exponentials need a precision check first.

\end{document}
